\documentclass[11pt]{article}

\usepackage[a4paper,margin=2.5cm]{geometry}
\usepackage{parskip}

\title{Difference Between PPO and PPO2}
\author{}
\date{}

\begin{document}

\maketitle

Compared to \texttt{ppo.py} in the maze 1 folder, \texttt{ppo2.py} uses more
guidance to solve a larger maze. The general PPO structure is the same, but it
adds BFS reward shaping, action masking, behavior cloning, and saving the best
model found during training.

The original reward function gives 100 points for reaching the exit, -1 point
for every normal step, -5 points for a wall hit, and an additional timeout
penalty of -50 points. In \texttt{ppo2.py}, BFS is used to calculate the
shortest distance from every position to the exit. A step toward the exit gives
0.17 points, while a step away from the exit gives -0.23 points. This tells the
model whether a move was useful.

The same eight features are used: the current position, the direction toward
the exit, and four flags describing whether the surrounding walls are blocked.
The BFS distance is only used to calculate rewards and is not given to the
model as an input feature.

\texttt{ppo2.py} also uses action masking. Moves that would hit a wall or leave
the maze are removed from the possible choices. This prevents the model from
wasting training steps on invalid moves.

Before PPO training begins, the model uses behavior cloning. BFS identifies
actions that move closer to the exit, and the model studies these correct
actions for 3000 epochs. Therefore, \texttt{ppo2.py} is not a pure
trial-and-error PPO model. It is behavior cloning followed by PPO fine-tuning.

The remaining PPO training is mostly unchanged. Each rollout gathers 2048
steps and is used for six training passes. Finally, \texttt{ppo2.py} remembers
the best model found during training and restores it before saving the weights.

\end{document}
